% =====================================================================================================
% Section IV. The Technology (sections/technology.tex)
% Sources: output/memos/m1_chinchilla.md (+_review), m2_techpanel.md (+_review), m8_measurement.md (+_review).
% Writer notes, number register and open issues: paper/notes/writer_technology.md.
% =====================================================================================================

\section{The Technology}\label{sec:tech}

Section~\ref{sec:ident} showed that curvature and the compute-optimal ratio are learned only from runs away from the
expansion path. Designed sweeps contain such runs, and this section uses them to estimate the parameter--data
technology: first on the canonical Chinchilla data (subsection A), then in seven public sweeps (B). We then ask whether better data are neutral or factor-biased (C), trace the
Kaplan--Chinchilla disagreement about the allocation exponent to input measurement and flexible inputs (D), and
describe our own controlled experiment (E). Three results emerge. Parameters and
data are gross complements in every sweep and under every functional form. The level of the elasticity depends on a
restriction that every sweep rejects, the unit outer exponent $\kappa=1$ of the Chinchilla form: $\sigma^*$ is
0.73--0.83 under that form and about 0.7 on the best-designed sweep once the data choose the curvature. And the objects
that allocation decisions need, the allocation exponent $a$ and the compute-optimal ratio $M^*$, are the least
portable across sweeps, estimators, and measurement conventions.

\subsection{The Canonical Data}\label{sec:tech:chinchilla}

We use Epoch's digitization of the runs in Figure~4 of \citet{hoffmann2022training}, the sample analyzed by
\citet{besiroglu2024chinchilla}: 245 runs on nine IsoFLOP budgets from $6\times10^{18}$ to $3\times10^{21}$ FLOP, with
tokens constructed as $D=C/(6N)$ (Section~\ref{sec:data}). Following \citet{besiroglu2024chinchilla}, the baseline
sample drops the five highest-loss runs ($n=240$). Figure~\ref{fig:chinchilla} shows the design in the input plane.
All curvature of the log reducible loss lies transverse to the expansion path (Section~\ref{sec:framework}); that
direction carries a variance of 0.90 in the design, against 1.83 along the path.

\begin{figure}[tp]
\centering
\includegraphics[width=\textwidth]{m1_chinchilla_isoquants}
\caption{The Chinchilla Design in the Input Plane}\label{fig:chinchilla}
\begin{figurenotes}
Epoch's digitization of the Chinchilla runs, shaded by validation loss (nats per token); hollow circles are the five
highest-loss runs dropped from the baseline sample. Curves: isoquants (iso-loss curves) of the Huber refit ($n=240$);
dashed lines: isocost lines $C=6ND$. Straight lines: expansion paths $N^*\propto C^{a}$ of the refit ($a=0.514$),
of Hoffmann et al.'s published Approach-3 parameters ($a=0.456$), and of Meta's Llama~3 Approach-2 law ($a=0.463$;
\citealt{grattafiori2024llama}). Arrows: design variance along the refit's path and transverse to it. Star:
Chinchilla-70B.
\end{figurenotes}
\end{figure}

% Table 3 -- Chinchilla data: estimator horse race, specification and duality tests, functional dependence.
% Condensed from output/tables/m1_chinchilla_horse_race.tex, m1_chinchilla_spec_tests.tex,
% m1_chinchilla_revealed_system.tex, m1_chinchilla_duality.tex, m1_chinchilla_identification.tex and the CSVs
% m1_chinchilla_{horse_race,estimator_diffs,system,spec_ces,spec_kappa,selection_rules,duality_classicalF,
% duality_bootcal,revealed_pref_70b,design_variance,fdep,fdep_sizecontrol}.csv; rank-one row from
% m2_spec_tests.csv (tau_q, p_rank_one_q). LAD column of the source table moved to the notes.
\begin{table}[tp]
\centering
\caption{The Chinchilla Data: Estimators and Specification Tests}
\label{tab:chinchilla}
\footnotesize
\setlength{\tabcolsep}{3pt}
\begin{tabular*}{\textwidth}{@{\extracolsep{\fill}}lcccc@{}}
\toprule
 & Huber & Gaussian & NLS, levels & System \\
 & (1) & (2) & (3) & (4) \\
\midrule
\multicolumn{5}{@{}l}{\textit{Panel A. Estimates} ($n=240$)} \\
$\alpha$ & 0.347 & 0.360 & 0.358 & 0.348 \\
 & (0.016) [0.024] & (0.013) [0.025] & (0.013) [0.026] & (0.014) \\
$\beta$ & 0.367 & 0.406 & 0.428 & 0.365 \\
 & (0.021) [0.046] & (0.025) [0.054] & (0.026) [0.055] & (0.016) \\
$a$ & 0.514 & 0.530 & 0.545 & 0.512 \\
 & (0.020) [0.044] & (0.017) [0.042] & (0.017) [0.042] & (0.015) \\
$\gamma$ & 0.178 & 0.191 & 0.195 & 0.178 \\
 & (0.006) [0.010] & (0.007) [0.014] & (0.007) [0.014] & (0.005) \\
$\sigma^*$ & 0.737 & 0.723 & 0.718 & 0.737 \\
 & (0.006) [0.010] & (0.008) [0.015] & (0.008) [0.015] & (0.005) \\
$M^*(5.76\times10^{23})$ & 17.9 & 14.3 & 10.4 & 18.5 \\
 & (6.9) [17.2] & (4.9) [12.1] & (3.8) [9.2] & (6.0) \\
$\hat\beta-\hat\beta_{(1)}$, paired & & 0.039 & 0.060 & \\
 & & (0.018) & (0.020) & \\
\addlinespace
\multicolumn{5}{@{}l}{\textit{Panel B. Specification tests} (Huber, $n=240$)} \\
 & \multicolumn{2}{c}{Estimate} & \multicolumn{2}{c}{$p$-value} \\
CES, $H_0$: $\alpha=\beta$ & \multicolumn{2}{c}{$\hat\alpha-\hat\beta=-0.020$ (0.029) [0.066]} & \multicolumn{2}{c}{0.50 [0.76]} \\
Outer exponent, $H_0$: $\kappa=1$ & \multicolumn{2}{c}{$\hat\kappa=0.774$ (0.060) [0.111]} & \multicolumn{2}{c}{$<$0.001 [0.042]} \\
\quad $\sigma^*$ with $\kappa$ free & \multicolumn{2}{c}{0.700 (0.015) [0.018]} & & \\
Rank-one curvature & \multicolumn{2}{c}{$\hat\tau=-0.009$} & \multicolumn{2}{c}{0.69} \\
System restrictions & \multicolumn{2}{c}{LR $=1.01$} & \multicolumn{2}{c}{0.60} \\
\addlinespace
\multicolumn{5}{@{}l}{\textit{Panel C. Duality and revealed preference, by technology}} \\
 & Refit (1) & Besiroglu & Hoffmann A3 & \\
Path test, classical $F_{2,7}$ ($p$) & 0.65 & 0.66 & 0.040 & \\
Path test, bootstrap-calibrated ($p$) & 0.23 & 0.18 & 0.060 & \\
Path level gap, $\Delta\ln N^*(\bar C)$ & 0.026 (0.046) & 0.028 (0.046) & 0.122 (0.046) & \\
$w$ at Chinchilla-70B & 1.04 [0.82, 1.42] & 1.03 & 0.62--0.71 & \\
\addlinespace
\multicolumn{5}{@{}l}{\textit{Panel D. Functional dependence} (outer exponent $\kappa$ free)} \\
 & Full design & On-path band & Random, same $n$ & \\
 & ($n=240$) & ($n=41$) & ($n=41$) & \\
Transverse variance share & 0.33 & 0.008 & 0.34 & \\
Condition number & 62 & 413 & 69--95 & \\
Max.\ profile LR for $\sigma^*$ & 424 & 2.0 & 79--101 & \\
\bottomrule
\end{tabular*}

\begin{tablenotes}
$L=E+AN^{-\alpha}+BD^{-\beta}$ on Epoch's digitization of the Chinchilla runs, five highest-loss runs dropped;
$a$, $\gamma$ and $\sigma^*$ as in Proposition~\ref{prop:duality} and equation~\eqref{eq:sigmaw}; $M^*(C)=D^*/N^*$ at
$C$ FLOP; $w=\varepsilon_N/\varepsilon_D$ at $N=70$B, $D=1.4$T. Estimators: (1) Huber loss ($\delta=10^{-3}$) on
$\ln L$; (2) Gaussian nonlinear least squares on $\ln L$; (3) nonlinear least squares on $L$; (4) system of the loss
equation (Laplace errors) and the nine IsoFLOP argmins, predicted design-consistently under the cross-equation
restrictions. Least absolute deviations on $\ln L$ gives $\alpha=0.348$, $\beta=0.366$, $\sigma^*=0.737$ and $w=1.03$.
Panel A: bootstrap standard errors (400 draws) resampling runs in parentheses and the nine IsoFLOP budgets in
brackets; for $M^*$, interquartile range/1.349; column 4, within-budget stratified bootstrap (300 draws).
Panel B: $\kappa$ from $L=E+[AN^{-a_1}+BD^{-b_1}]^{\kappa}$, $\sigma^*=2/(2+a_1+b_1)$;
rank one: $\hat\tau=\det H/\lVert H\rVert_F^2$ for a translog of $\ln(L-\hat E)$, $\hat E$ from the $\kappa$ family,
residual-bootstrap null; system: nested LR of a level and slope shift of the argmin equation, $\chi^2_2$. Panel C:
design-consistent tests of the expansion-path restrictions (the Approach-2 procedure applied to each technology's
fitted values at the observed design); bootstrap-calibrated $p$ from the stratified bootstrap distribution of the Wald
statistic (wild bootstrap: 0.70, 0.75, 0.21); level gap: argmins minus null at the geometric-mean budget; $w$ interval:
pairs bootstrap (budget bootstrap: [0.54, 2.01]); the estimators in columns 2--4 give $w=1.14$, 1.30 and 1.03; Hoffmann A3 at rounded and full published precision; further tests in Online Appendix
Table~\ref{tab:app-duality}. Panel D: on-path band
$|\ln N-\ln N^*(C)|\le0.15$; random: four subsamples of the full design of the band's size; condition number of the
Jacobian in normalized parameters; profile LR over $\sigma^*\in[0.50,0.95]$.
\end{tablenotes}

\begin{tablenotes}[Source]
\citet{hoffmann2022training}; \citet{besiroglu2024chinchilla}; authors' calculations.
\end{tablenotes}
\end{table}


Table~\ref{tab:chinchilla}, panel A, reports four estimators on the same runs. The reference estimator of
\citet{hoffmann2022training} and \citet{besiroglu2024chinchilla}, a Huber loss with threshold $\delta=10^{-3}$ on
$\ln L$, is least absolute deviations (LAD) in practice: 84 percent of its residuals lie in the linear region, and
across bootstrap draws the Huber and LAD estimates of $\beta$ differ with a standard deviation of 0.004. Besiroglu et
al.'s published parameters are in fact an LAD-type optimum, a Huber likelihood whose estimated scale places every
residual in the linear region; our LAD estimate reproduces them to four or five digits. The two robust estimators
give $\beta=0.367$, Gaussian nonlinear least squares on $\ln L$ gives 0.406, and least squares in levels 0.428.
Because all estimators use the same bootstrap draws, the differences are precisely estimated: 0.039 (s.d.\ 0.018)
and 0.060 (0.020). We prefer the robust estimators because the residuals are heavy-tailed, as shown below. The
disagreement matters less for the derived objects: $\sigma^*$ ranges from 0.718 to 0.737 and $a$ from 0.513 to 0.545,
and the implied differences in $M^*$ at Chinchilla's compute, $-3.7$ (4.9) and $-7.6$ (5.3), are not significant.

Runs on the same budget share a design and plausibly common shocks. Resampling the nine budgets rather than runs
raises the standard errors of the exponents by a factor of 1.5 to 2.2 (for $\beta$, from 0.021 to 0.046) and widens the 95 percent
interval for $M^*(5.76\times10^{23})$ from [7.8, 35.1] to [3.2, 126]. With nine clusters the cluster bootstrap is
itself imprecise \citep{cameron2008bootstrap}, but the conclusion is robust: these data alone do not determine the
compute-optimal ratio at frontier scale.

Besiroglu et al.'s exclusion of the five highest-loss runs is, formally, truncation on the dependent variable, and it
matters: $\hat\beta$ falls from 0.453 to 0.367 (paired difference 0.086, s.e.\ 0.052) and keeps falling, almost
monotonically, as more runs are dropped (0.347 at fifteen). The mechanism is not truncation bias. A truncated-Gaussian
estimator \citep{hausman1977social} on the retained runs returns the untruncated estimate ($\hat\beta=0.406$ in both),
because no retained run lies near the threshold. The five runs are instead the most data-starved runs of the
$10^{19}$ budget ($D/N$ between 0.04 and 0.40), with log residuals 9--45 times the residual standard deviation;
dropping runs with $D/N<0.4$, selection on the regressor, gives $\hat\beta=0.371$. The swing is the leverage of a few
gross outliers where the functional form fails as $D/N\to0$. (The truncated estimator maintains Gaussian errors, which
these outliers contradict.) The exclusion is consequential for revealed preference below: the wedge of Chinchilla-70B
is 1.54 with the outliers and between 0.95 and 1.06 under every rule that removes them.

Proposition~\ref{prop:duality} reads Chinchilla's three approaches as estimation of the cost function (Approach 1),
the conditional factor demands (Approach 2), and the primal technology (Approach 3), linked by cross-equation
restrictions. Approach 2 on the extraction gives $a=0.498$ (s.e.\ 0.025), against the 0.49 that
\citet{hoffmann2022training} report. Testing the restrictions requires care, because parabola argmins on a finite grid
are biased \citep{czech2026problems}: the null must be the Approach-2 procedure applied to the primal's fitted values
at the observed design. Naive tests against the analytic path reject every technology, including our refit.
Design-consistent tests reject neither our refit nor Besiroglu et al.'s parameters ($p\ge0.16$ in every classical or
bootstrap-calibrated variant; panel C and Online Appendix Table~\ref{tab:app-duality}). Hoffmann et al.'s published Approach-3 parameters are marginally inconsistent
with their own IsoFLOP argmins: $p=0.040$ in a classical $F$ test on the nine argmin gaps, 0.060 with
bootstrap-calibrated critical values, and 0.21--0.38 under a wild bootstrap. The inconsistency is in the level of the path.
Those parameters imply a compute-optimal model about 12 percent smaller than the argmins show ($\Delta\ln N^*=0.122$,
s.e.\ 0.046); the slope gap (0.042, s.e.\ 0.023) is not significant. The cost-function restrictions fare worse for
every technology: the Approach-1 minima lie 0.005--0.008 nats below each Approach-3 frontier, a rejection under the
design-resampling bootstrap that does not survive the wild bootstrap.

Duality also yields a revealed-preference test. If Chinchilla-70B ($N=70$B, $D=1.4$T) minimized training compute,
its wedge $w=\varepsilon_N/\varepsilon_D$ equals one (Proposition~\ref{prop:wedge} with $T=0$). Our refit gives
$w=1.04$ (95 percent interval [0.82, 1.42]). Hoffmann et al.'s parameters imply 0.62 at their rounded values and 0.71
at full precision, outside that interval but inside the nine-cluster interval [0.54, 2.01]. Because the statement is joint,
about the parameters and the optimality of the chosen model, we read it as consistency of the refit with the choice,
not as a rejection of Hoffmann et al.'s parameters. A system estimator in the spirit of
\citet{leonledesma2010identifying}, which combines the loss equation with the design-consistent argmin equation,
imposes the restrictions directly (column 4). The data do not reject them (nested LR 1.01, $p=0.60$), and the system
estimates are those of the primal with a modest gain in precision (s.e.\ of $a$ 0.015, against 0.018 under the same
bootstrap).

Panel B of Table~\ref{tab:chinchilla} reports specification tests. The homothetic CES restriction $\alpha=\beta$ is not rejected. The outer-exponent
restriction is. Nesting Chinchilla in the family $L=E+[AN^{-a_1}+BD^{-b_1}]^{\kappa}$ of
equation~\eqref{eq:ident-family}, which contains Kaplan's form at $\kappa\approx0.1$, gives $\hat\kappa=0.774$ (s.e.\
0.060; 0.111 with budget clusters), so $\kappa=1$ is rejected ($p<0.001$; cluster $p=0.042$). The restriction is not
innocuous for curvature: with $\kappa$ free, $\sigma^*=2/(2+a_1+b_1)$ falls from 0.737 to 0.700 (0.015), while the
on-path objects barely move ($a=0.504$, $\gamma=0.165$). Because $\kappa$ is identified from the shape of digitized
IsoFLOP profiles, whose losses are quantized, we rely on the replications below. The rank-one curvature implied by
every member of the family (Online Appendix~A, Lemma~\ref{lem:geometry}) is not rejected once $E$ is estimated in the $\kappa$ family
($p=0.69$), although translog tests of it are sensitive to $\hat E$.

Finally, the data display the functional dependence of Proposition~\ref{prop:fdep}. The 41 runs within 0.15 log
points of the fitted expansion path have a transverse variance share of 0.8 percent (panel D). With $\kappa$ free, the
profile likelihood over $\sigma^*$ is flat on this subsample: the likelihood-ratio statistic never exceeds 2.0 on
$[0.50,0.95]$, so every value lies in the 95 percent confidence set. On the full design it reaches 424, and the set is
[0.67, 0.69]. The flatness is a design effect, not a sample-size effect: random subsamples of 41 runs from the full
design reach 79--101. Relative to those subsamples, on-path standard errors are 2.3 times larger for $\alpha$ and 2.2
times for $a$, but not larger for the frontier elasticity $\gamma$ (0.94): the frontier is identified on the path;
the curvature is not. Under $\kappa=1$, on-path $\sigma^*$ remains precise (s.e.\ 0.011), because the functional form,
not the data, pins it down.

\subsection{Across Sweeps}\label{sec:tech:sweeps}

Table~\ref{tab:technology} estimates the technology in the public sweeps described in Section~\ref{sec:data}:
Chinchilla, Farseer \citep{li2025predictableb}, the three web corpora of \citet{gadre2024language}, the OLMo ladder
\citep{bhagia2024establishing}, the single-epoch runs of \citet{muennighoff2023scaling}, and DataDecide
\citep{magnusson2025datadecide}. The designs differ. The dispersion of $\ln M$ given $\ln C$, the transverse variation
that identifies curvature (Proposition~\ref{prop:info}), ranges from 1.97 (Farseer, Muennighoff et al.) to 0.73
(DataDecide, whose final checkpoints all have $D\approx100N$ and whose identification comes from checkpoints taken
before the learning-rate schedule ends). Loss units and tokenizers also differ, so only the unit-free objects are
compared across rows.

\begin{figure}[tp]
\centering
\includegraphics[width=\textwidth]{m2_fig3_sigma_forest}
\caption{The Elasticity of Substitution across Sweeps}\label{fig:sigma}
\begin{figurenotes}
Panel a: estimates with 95 percent bootstrap intervals by sweep. Circles and squares: $\sigma^*=2/(2+\alpha+\beta)$
under the Chinchilla form, estimated by Huber loss and by Gaussian nonlinear least squares on $\ln L$; diamonds:
$\sigma^*_\kappa=2/(2+a_1+b_1)$ with the outer exponent free (labeled $q$ in the legend); triangle: median local
elasticity from a kernel-weighted local quadratic (Farseer). Open symbols: DataDecide, pooled fit with common exponents,
identified from checkpoints before the end of the learning-rate schedule. Shaded band: the range of capital--labor
estimates, 0.4--0.7. Panel b: $\sigma^*$ from separate Chinchilla-form fits to each of the 25 DataDecide recipes
(Huber); dashed line: pooled estimate.
\end{figurenotes}
\end{figure}

% Table 4 -- The parameter-data technology across public sweeps.
% Condensed from output/tables/m2_table3_technology.tex / .csv (Panel A, Huber rows; kappa columns = q columns)
% and m2_spec_tests.csv; Panel B from m2_farseer_forms.csv, m2_farseer_local_sigma_summary.csv,
% m2_table3_technology.csv (sigma_star_q CI). Notation: the module's outer exponent q is written kappa here.
\begin{table}[tp]
\centering
\caption{The Parameter--Data Technology across Public Training Sweeps}
\label{tab:technology}
\footnotesize
\setlength{\tabcolsep}{1.5pt}
\begin{tabular*}{\textwidth}{@{\extracolsep{\fill}}lcccccccccc@{}}
\toprule
\multicolumn{11}{@{}l}{\textit{Panel A. Estimates by sweep} (Huber loss on $\ln L$)} \\
 & & & \multicolumn{6}{c}{Chinchilla form ($\kappa=1$)} & \multicolumn{2}{c}{$\kappa$ free} \\
\cmidrule(lr){4-9}\cmidrule(l){10-11}
Sweep & Runs & $s_M$ & $\alpha$ & $\beta$ & $a$ & $\gamma$ & $\sigma^*$ & $M^*(10^{21})$ & $\hat\kappa$ & $\sigma^*_\kappa$ \\
\midrule
Chinchilla & 240 & 1.35 & 0.347 & 0.367 & 0.514 & 0.178 & 0.737 & 21.4 & 0.77 & 0.701 \\
 & & & (0.015) & (0.021) & (0.020) & (0.006) & (0.006) & [15.1, 28.6] & (0.06) & (0.014) \\
Farseer & 404 & 1.97 & 0.281 & 0.312 & 0.526 & 0.148 & 0.772 & 24.8 & 0.47 & 0.710 \\
 & & & (0.028) & (0.011) & (0.022) & (0.009) & (0.010) & [18.8, 32.9] & (0.01) & (0.003) \\
Gadre, C4 & 34 & 1.46 & 0.253 & 0.262 & 0.509 & 0.129 & 0.795 & 6.9 & 0.35 & 0.607 \\
 & & & (0.046) & (0.056) & (0.072) & (0.015) & (0.020) & [1.3, 88] & (0.08) & (0.054) \\
Gadre, RedPajama & 35 & 1.47 & 0.243 & 0.273 & 0.529 & 0.129 & 0.795 & 5.8 & 0.39 & 0.623 \\
 & & & (0.027) & (0.042) & (0.045) & (0.009) & (0.013) & [1.2, 13.3] & (0.14) & (0.059) \\
Gadre, RefinedWeb & 35 & 1.47 & 0.223 & 0.290 & 0.566 & 0.126 & 0.796 & 3.4 & 0.36 & 0.594 \\
 & & & (0.035) & (0.038) & (0.055) & (0.012) & (0.015) & [1.2, 18.9] & (0.12) & (0.062) \\
OLMo ladder & 30 & 0.94 & 0.267 & 0.157 & 0.370 & 0.099 & 0.825 & 48 & 0.22 & 0.544 \\
 & & & (0.118) & (0.041) & (0.177) & (0.022) & (0.034) & [5.9, 298] & (0.02) & (0.021) \\
Muennighoff & 33 & 1.97 & 0.336 & 0.387 & 0.535 & 0.180 & 0.735 & 60 & 0.42 & 0.511 \\
 & & & (0.131) & (0.077) & (0.126) & (0.041) & (0.031) & [9.3, 281] & (0.10) & (0.056) \\
DataDecide & 1,100 & 0.73 & 0.284 & 0.131 & 0.315 & 0.090 & 0.828 & -- & & \\
 & & & (0.023) & (0.013) & (0.037) & (0.004) & (0.005) & & & \\
\bottomrule
\end{tabular*}

\vspace{6pt}
\begin{tabular*}{\textwidth}{@{\extracolsep{\fill}}lccccc@{}}
\multicolumn{6}{@{}l}{\textit{Panel B. Farseer: functional forms and the elasticity without $\kappa=1$}} \\
\midrule
 & & \multicolumn{3}{c}{RMSE of $\ln L$} & \\
\cmidrule(lr){3-5}
Specification (elasticity) & Estimate & In sample & Top $C$ decile & Largest $N$ & BIC \\
\midrule
Chinchilla form ($\sigma^*$) & 0.772 (0.010) & 0.0087 & 0.025 & 0.013 & $-2{,}647$ \\
$\kappa$ family ($\sigma^*_\kappa$) & 0.710 [0.705, 0.716] & 0.0025 & 0.0052 & 0.0036 & $-3{,}663$ \\
Farseer form (median $\sigma$) & 0.706 & 0.0018 & 0.0026 & 0.0026 & $-3{,}916$ \\
Local quadratic (median $\sigma$) & 0.690 [0.679, 0.699] & & & & \\
\bottomrule
\end{tabular*}

\begin{tablenotes}
Panel A: $L=E+AN^{-\alpha}+BD^{-\beta}$ estimated separately for each sweep by Huber loss ($\delta=10^{-3}$) on
$\ln L$. $s_M$: residual standard deviation of $\ln(D/N)$ given $\ln C$ (off-path design variation).
$a=\beta/(\alpha+\beta)$; $\gamma=\alpha\beta/(\alpha+\beta)$; $\sigma^*=2/(2+\alpha+\beta)$; $M^*(10^{21})$:
compute-optimal $D/N$ at $C=6ND=10^{21}$ FLOP (extrapolated where the largest budget is smaller), with 95 percent
percentile interval. Last two columns: $L=E+[AN^{-a_1}+BD^{-b_1}]^{\kappa}$, $\sigma^*_\kappa=2/(2+a_1+b_1)$;
$\kappa=1$ is rejected in every sweep ($p<0.001$, bootstrap Wald and Gaussian LR). Bootstrap standard errors in
parentheses: 400 draws resampling runs (Chinchilla, Farseer, Gadre, Muennighoff) or (size, multiplier) cells (OLMo
ladder); 200 draws for the $\kappa$ columns. Chinchilla standard errors differ slightly from Table~\ref{tab:chinchilla},
which uses separate bootstrap draws. Gadre: three web corpora of \citet{gadre2024language}; Muennighoff: single-epoch runs. DataDecide: 25 recipes pooled with common exponents and recipe-specific
$(A_r,B_r,E_r)$, estimated on 21,888 checkpoints of 1,100 runs taken before the learning-rate schedule ends, 100
draws resampling (size, seed) cells; levels are descriptive and $M^*$ is not reported. Loss units and tokenizers
differ across rows (Farseer: bits per character); $\alpha$, $\beta$, $a$, $\gamma$, $\sigma$, and $\kappa$ are unit
free. Parameter counts: non-embedding (Farseer), total (Chinchilla, Gadre), excluding the input embedding only (OLMo
ladder, DataDecide). Gaussian least squares, a pooled Gadre fit, and robustness checks are in Online Appendix Tables~\ref{tab:app-robust} and~\ref{tab:app-farseer}.
Panel B: Chinchilla form by Huber loss; $\kappa$ family and Farseer's nine-parameter form (Eq.~3 of
\citealt{li2025predictableb}) by Gaussian least squares, except $\sigma^*_\kappa$ (Huber, as in Panel A). Held-out
errors: models refitted without the top decile of $6ND$ or without the four largest model sizes, from
training-sample starting values. Local quadratic: Gaussian-kernel local quadratic of $\ln L$ in $(\ln N,\ln D)$ at 42
interior points, with a bootstrap interval for the median; $\sigma$ in the Farseer form is evaluated at the same
points.
\end{tablenotes}

\begin{tablenotes}[Source]
\citet{hoffmann2022training}; \citet{besiroglu2024chinchilla}; \citet{li2025predictableb}; \citet{gadre2024language};
\citet{bhagia2024establishing}; \citet{muennighoff2023scaling}; \citet{magnusson2025datadecide}; authors' calculations.
\end{tablenotes}
\end{table}


Parameters and data are gross complements in every sweep. Under the Chinchilla form $\sigma^*$ ranges from 0.735
(Muennighoff et al.) to 0.828 (DataDecide), with standard errors between 0.005 and 0.034, and no interval approaches
one (Figure~\ref{fig:sigma}). Lemma~\ref{lemma:sigma} requires $\sigma<1$ of any technology with interior compute
optima, so this is a check on the data rather than a discovery; the economic content is the level.

The level depends on the outer exponent. The restriction $\kappa=1$ is rejected in all seven sweeps ($p<0.001$;
$\hat\kappa$ between 0.22 and 0.77), and relaxing it improves the fit substantially (the root mean squared error of
$\ln L$ on Farseer falls from 0.0085 to 0.0025). With $\kappa$ free, $\sigma^*_\kappa$ ranges from 0.51 to 0.71. In
the 30--35-run sweeps the inner parameters are weakly identified and the intervals are about $\pm0.1$ wide. Farseer,
with 404 runs on a near-factorial design, pins the level down: $\sigma^*_\kappa=0.710$ [0.705, 0.716]. Two estimators
that do not impose the Chinchilla form agree (Table~\ref{tab:technology}, panel B). A kernel-weighted local quadratic of $\ln L$ in
$(\ln N,\ln D)$, which needs no estimate of $E$ because $\sigma$ is ordinal, gives a median local elasticity of 0.690
[0.679, 0.699] over 42 interior points, and Farseer's own functional form implies a median of 0.706. (Local estimates
are noisy near the edge of the design, from 0.61 to 0.99, so we use only the median.) The Chinchilla-implied 0.77
lies outside both intervals. The Chinchilla form also extrapolates poorly in compute: fitted without the top decile of
compute, it over-predicts the held-out losses, with an out-of-sample error of 0.022--0.025 against 0.0026 for
Farseer's form.

We therefore report $\sigma$ as a range: 0.73--0.83 under the Chinchilla restriction and about 0.7 when the data
choose the curvature. Estimates of the capital--labor elasticity concentrate between 0.4 and 0.7
\citep{chirinko2008sigma,klump2007factor,oberfield2021micro}, with micro estimates lower still
\citep{raval2019micro}. Under the Chinchilla form, parameters and data sit just above that range; once the curvature
is freed, the estimates overlap its top. Parameters and data are complements of roughly the strength of capital and
labor, if anything slightly weaker.

The other restrictions of the Chinchilla form survive. CES is not rejected by a Wald test in any sweep ($p$ between
0.22 and 0.90). Rank-one curvature, with $E$ from the $\kappa$ family, is not rejected at 5 percent in six of seven
sweeps; on Farseer a small departure is detectable ($p\le0.003$), but its bootstrap interval contains zero. The data
give little evidence against the additive separability behind the geometry of Section~\ref{sec:ident}.

What differs across sweeps is the allocation. The allocation exponent ranges from 0.37 (OLMo ladder) to 0.57 (Gadre
et al., RefinedWeb), and $M^*(10^{21})$ from 3.4 to 60, with single-sweep intervals spanning factors of 1.7 to 70;
$\sigma^*$ (0.735--0.825) and the frontier elasticity $\gamma$ (0.10--0.18) vary by less than a factor of two. Within a
sweep, $a$ and $M^*$ also respond to choices that barely move $\sigma^*$. Counting Farseer's untied embeddings moves
$a$ from 0.526 to 0.411 and $M^*(10^{23})$ from 19.5 to 45 ($\sigma^*$ from 0.772 to 0.724). Dropping Muennighoff et
al.'s two runs with $D/N<0.4$ moves $a$ from 0.535 to 0.360 and $M^*(10^{21})$ from 60 to 139 ($\sigma^*$ from 0.735
to 0.726). On Gadre et al.'s RefinedWeb runs, $M^*(10^{21})$ is 3.4 under Huber loss and 10.4 under least squares. We
do not read these contrasts as evidence that $\sigma$ is intrinsically better determined than $a$. The map
$\sigma^*=2/(2+\alpha+\beta)$ compresses variation in $\alpha+\beta$ by a factor of $\sigma^{*2}/2\approx0.3$, and on
the Step Law grid of subsection D $\sigma^*$ is no more stable than $a$ in units of its own sampling error. The
economically relevant point is Proposition~\ref{prop:fdep}: $M^*$ depends on the level ratio $A/B$, which designed
variation identifies only weakly and which absorbs differences in units, conventions, and recipes. The frontier
elasticity and $\sigma$ are comparable across sweeps in point estimate; $M^*$ is not.

\subsection{Data Quality: Neutral or Factor-Biased?}\label{sec:tech:quality}

Is better data a neutral improvement, a data-augmenting one, or only a lower asymptote? The answer matters for
allocation and for the wedge of Section~\ref{sec:wedge}. At given compute, Hicks-neutral productivity does not move
the optimal input mix, data-augmenting productivity lowers the optimal tokens per parameter, and factor bias
$\chi=\beta\psi_D-\alpha\psi_N$, where $\psi_N$ and $\psi_D$ are the factor-augmenting productivity terms of
Section~\ref{sec:framework}, contaminates the revealed wedge, $\hat w=we^{-\chi}$
(Proposition~\ref{prop:wedge}). We nest these hypotheses in $L_r=E_r+A_rN^{-\alpha}+B_rD^{-\beta}$ for corpus or
recipe $r$, evaluate every model on the same C4 validation set so that output is common, and base inference on a wild
cluster bootstrap with a Wald cross-check (Online Appendix Table~\ref{tab:app-neutrality}).

Across Gadre et al.'s three web corpora (104 runs), the corpora differ ($p\le0.002$ against a single technology), and
an asymptote shift alone is rejected ($p=0.004$) although it accounts for 90 percent of the cross-corpus variation in
fit. Hicks-neutral scaling of reducible loss with corpus-specific asymptotes is not rejected against common exponents
($p=0.90$); it accounts for 96.8 percent, against 96.9 percent for unrestricted common exponents. Power is low:
data- and parameter-augmenting alternatives with free asymptotes are not rejected by the wild bootstrap either ($p=0.13$
and 0.15), and a Wald test rejects the data-augmenting version under least squares but not under Huber loss. Whatever
the bias, it is small: compute-optimal $M$ differs by 8--18 percent across corpora.

DataDecide's 25 recipes (21,888 checkpoints) have the power Gadre et al.'s three corpora lack, and every restriction
is rejected; Hicks neutrality with free asymptotes gives an LR of 145 on 24 degrees of freedom (bootstrap $p\le0.01$).
Yet the economic content is concentrated in the asymptote and a neutral scale. Recipe-specific asymptotes account for
98.1 percent of the cross-recipe variation in fit, a neutral scale raises this to 99.71 percent, and factor bias adds
0.05 points. The biased component is small in loss but systematic, and it matters for allocation: recipe tilts
$\ln(A_r/B_r)$ span 0.22--0.26, which implies compute-optimal $M$ varying 2.9--3.4-fold across recipes and a revealed
wedge varying 1.25--1.29-fold at given inputs. The data-augmenting tilt is carried by the heavily filtered DCLM
variants, whose C4 asymptotes are the highest of all recipes. Whether ``better'' data are data-augmenting is
therefore not identified here; under common exponents a lower $M^*$ given the tilt is an identity. Moreover, only
relative tilts are interpretable, and only if the unfinished schedule affects all recipes alike, and asymptote shifts
on a common validation set partly measure distribution match rather than quality. The elasticity itself is homogeneous across recipes:
$\sigma^*_r$ lies between 0.790 and 0.839 (Figure~\ref{fig:sigma}, panel b).

In the public sweeps, then, better data mostly lower the asymptote and scale reducible loss neutrally. The factor-biased
part is small in loss but large enough to move $M^*$ threefold and the wedge by 25--29 percent, which
Section~\ref{sec:wedge} carries as a rival explanation for over-training.

\subsection{Measurement and Flexible Inputs}\label{sec:tech:measurement}

The largest disagreement in the scaling literature concerns the allocation exponent: $N^*\propto C^{0.73}$ in
\citet{kaplan2020scaling} against about 0.5 in \citet{hoffmann2022training}. \citet{porian2024resolving} trace it to a
sequence of changes in training practice, and \citet{pearce2024reconciling} to the treatment of embedding parameters.
Both are classic problems in production-function estimation. Counting the parameters that cost FLOPs, such as the
output head, is input measurement, whose bias depends on how the measurement error
correlates with scale \citep{nerlove1963returns,christensen1976economies,collardwexler2016production}. Warmup,
learning-rate schedule, learning rate, and batch size are flexible inputs; when they are not re-optimized at each
scale, they are not concentrated out of the production function \citep{ackerberg2015identification,gandhi2020identification}.

% Table 5 -- Measurement and flexible inputs in the allocation exponent.
% Condensed from output/tables/m8_measurement_porian.tex (Panels A-B; CSVs m8_measurement_porian_steps.csv,
% m8_measurement_porian_counting.csv). The former Panel C (Chinchilla non-embedding, Step Law SFA and demand) moved to
% the text and Online Appendix Tables D17-D18 at integration (2026-09-24).
\begin{table}[tp]
\centering
\caption{Measurement and Flexible Inputs in the Allocation Exponent}
\label{tab:measurement}
\footnotesize
\setlength{\tabcolsep}{2pt}
\begin{tabular*}{\textwidth}{@{\extracolsep{\fill}}lcccccc@{}}
\toprule
\multicolumn{7}{@{}l}{\textit{Panel A. Porian et al.'s path, re-estimated from the 975 released runs}} \\
 & \multicolumn{3}{c}{RefinedWeb} & \multicolumn{3}{c}{OpenWebText2} \\
\cmidrule(lr){2-4}\cmidrule(l){5-7}
Step & $\hat a$ & Published & $\Delta\hat a$ & $\hat a$ & Published & $\Delta\hat a$ \\
\midrule
1. Kaplan reproduction & 0.837 & 0.835 & & 0.862 & 0.864 & \\
 & [0.823, 0.845] & & & [0.817, 0.896] & & \\
2. + head in $N$ and $C$ & 0.705 & 0.706 & $-0.132$ & 0.700 & 0.699 & $-0.162$ \\
 & [0.691, 0.716] & & & [0.666, 0.722] & & \\
3. + warmup $\propto N$ & 0.602 & 0.602 & $-0.103$ & 0.602 & 0.603 & $-0.098$ \\
 & [0.591, 0.617] & & & [0.570, 0.631] & & \\
4. + cosine decay & 0.572 & 0.571 & $-0.030$ & 0.574 & 0.574 & $-0.028$ \\
 & [0.561, 0.588] & & & [0.543, 0.604] & & \\
5. + tuned LR, batch, $\beta_2$ & 0.496 & 0.497 & $-0.075$ & 0.517 & 0.518 & $-0.057$ \\
 & [0.490, 0.504] & & & [0.491, 0.541] & & \\
Share of step 2 in total & & & 0.39 & & & 0.47 \\
\addlinespace
\multicolumn{7}{@{}l}{\textit{Panel B. Same runs, only the parameter count changes (training loss)}} \\
 & Kaplan setup & Tuned & Diff. & Kaplan setup & Tuned & Diff. \\
\cmidrule(lr){2-4}\cmidrule(l){5-7}
Non-embedding (Kaplan) & 0.837 & 0.620 & 0.217 & 0.862 & 0.631 & 0.231 \\
 & (0.006) & (0.005) & & (0.020) & (0.017) & \\
+ head (standard) & 0.700 & 0.502 & 0.198 & 0.696 & 0.501 & 0.194 \\
 & (0.006) & (0.004) & & (0.017) & (0.012) & \\
+ head, attention FLOPs & 0.717 & 0.505 & 0.212 & 0.713 & 0.507 & 0.206 \\
 & (0.006) & (0.004) & & (0.018) & (0.013) & \\
+ head, input embeddings & 0.661 & 0.460 & 0.202 & 0.670 & 0.471 & 0.199 \\
 & (0.007) & (0.003) & & (0.017) & (0.012) & \\
\bottomrule
\end{tabular*}


\begin{tablenotes}
$a$ is the exponent of $N^*\propto C^{a}$ from the IsoFLOP procedure of \citet{porian2024resolving}, re-implemented
on their released runs (16 architectures, 5M--901M parameters; budgets $1.25\times10^{16}\cdot2^i$,
$i=0,\dots,11$). Panel A: step 2 is input measurement, steps 3--5 flexible inputs; brackets, 95 percent intervals from their noise-and-interpolate bootstrap (1,000 draws);
``Published'' = their Table~1; step 1 uses Kaplan's count and training loss, steps 2--5 the standard count and
validation loss. Panel B: ``Kaplan setup'' = untuned runs with Kaplan's 1.57B-token warmup; ``Tuned'' = per-size tuned
learning rate, batch, and $\beta_2$ with a constant learning rate; bootstrap standard errors in parentheses. Evidence from the Chinchilla data with non-embedding $N$ and from the Step Law grid (stochastic frontier,
learning-rate demand) is reported in the text and in Online Appendix Tables~\ref{tab:app-sfa} and~\ref{tab:app-demand}.
\end{tablenotes}

\begin{tablenotes}[Source]
\citet{porian2024resolving}; \citet{hoffmann2022training}; \citet{besiroglu2024chinchilla}; \citet{li2025predictablea};
\citet{bjorck2024scaling}; authors' calculations.
\end{tablenotes}
\end{table}


Re-estimating Porian et al.'s path from their 975 released runs reproduces all ten published exponents to within
0.003 (Table~\ref{tab:measurement}, panel A). The exponent falls from 0.837 to 0.496 on RefinedWeb and from 0.862 to
0.517 on OpenWebText2. Counting the head accounts for 39 and 47 percent of the drop; flexible inputs account for the
remaining 53--61 percent. Because a sequential decomposition depends on the order of the steps, panel B holds the runs
fixed and changes only the count. Counting the head lowers $a$ by 0.12--0.17 whether or not the runs are tuned, and
tuning lowers it by 0.19--0.23 under every count, so the two components are nearly additive. These models have 5M to
901M parameters.

The measurement component has a closed form. Let the true parameter count be $N=N_m+X(N_m)$, where the omitted
component $X$ has share $s=X/N$ and scale elasticity $\theta=d\ln X/d\ln N_m<1$. The IsoFLOP argmin in measured units
then has local slope
\begin{equation}\label{eq:measbias}
a_m=\frac{\beta}{\beta+\alpha\varphi-\zeta},\qquad \varphi\equiv1-s(1-\theta),\qquad
\zeta\equiv\frac{d\ln\varphi}{d\ln N_m}=\frac{(1-\theta)s(\varphi-\theta)}{\varphi}\ge0,
\end{equation}
so $a_m\ge a$, with equality when $s=0$ or $\theta=1$ (derivation in Online Appendix~\ref{app:add-biasformulas}). A proportional error moves
only the level of the path. The bias comes from the scale-correlated part of the error, as in
\citet{collardwexler2016production}, and it vanishes at scale as the omitted share falls, which is Nerlove's lesson
that a local elasticity estimated over a restricted size range is not structural. With Pearce and Song's embedding
map, equation~\eqref{eq:measbias} predicts a measured exponent over Kaplan's size range of 0.792 under Besiroglu et
al.'s parameters and 0.757 under Hoffmann et al.'s, against the 0.78 and 0.74 that Pearce and Song obtain by
simulation. On Porian et al.'s ladder it gets the sign and magnitude roughly right but over-predicts the bias, by 20--50
percent on the actual design. The Chinchilla data are not immune: with non-embedding parameter counts the allocation
exponent rises from 0.514 to 0.556 (paired difference 0.042, s.e.\ 0.006).

Flexible-input inefficiency has an analogous first-order formula. Write observed loss as $L^*e^{\iota(N,D)}$, where
$L^*$ is the concentrated technology and $\iota\ge0$ the inefficiency of the hyperparameter policy. Along an IsoFLOP
curve the log argmin $\ln N^*$ shifts by $-\Delta(C)/f''(C)$, where $\Delta$ is the slope of the excess loss in $\ln N$ at fixed
compute and $f''=(\alpha+\beta)\gamma R^*$ is the curvature of the concentrated IsoFLOP curve, with $R^*$ the reducible
loss on the frontier. Hence $a_{\rm obs}-a=-\partial[\Delta/f'']/\partial\ln C$. Only the transverse gradient of the
inefficiency matters: inefficiency that depends on the inputs only through compute leaves $a$ unbiased. The bias in
$a$ is first order in that gradient, although, by the envelope theorem, the effect on the loss level is second order
in the hyperparameter error. On Porian et al.'s runs the formula explains about half to all (47--114 percent,
depending on how $f''$ is estimated) of the three flexible-input steps. The mechanism is Kaplan et al.'s fixed
1.57-billion-token warmup and large batch, which leave short runs with too few useful steps. At $2.5\times10^{16}$
FLOP the untuned runs' excess loss at the tuned optimum is 1.57 nats; it vanishes by about $3\times10^{18}$ FLOP. A
penalty of this kind pushes small-budget argmins toward smaller models and steepens the fitted path.

The Step Law grid \citep{li2025predictablea}, with 1,911 runs in 17 $(N,D)$ cells and up to 120 learning-rate and
batch-size configurations per cell, shows the flexible inputs directly (Online Appendix Tables~\ref{tab:app-steplaw}--\ref{tab:app-demand}). A stochastic frontier with
heteroskedastic inefficiency \citep{aigner1977formulation,caudill1995frontier} shows that inefficiency is data-biased:
its scale falls with $D$ (elasticity $-0.23$, cluster-robust s.e.\ 0.03) but not detectably with $N$ ($-0.04$, s.e.\
0.08); longer training flattens the loss landscape in the hyperparameters. The estimated demand functions replicate
Step Law's learning-rate rule, and the distinction between conditional and unconditional demand
\citep{milgrom1996lechatelier} explains part of a sign disagreement in the literature: Step Law's positive elasticity
of the optimal learning rate in $D$ lets the batch adjust, whereas the $-0.32$ of \citet{bjorck2024scaling} holds the
batch fixed. Holding the batch fixed in the Step Law data lowers the elasticity to 0.046 (0.015), removing the
positive sign, but this closes only 26--43 percent of the gap between the two published values. The Step Law design
itself, with five model sizes and no IsoFLOP profiles, does not identify $a$: the sum of squared residuals varies by
less than 17 percent for $E$ between 0.6 and 1.7, over which $a$ runs from 0.76 to 0.45 and $\sigma^*$ moves by 7.0 of
its own standard errors (against 2.8 for $a$).

The allocation exponent is thus the object most exposed to measurement conventions and to flexible inputs that are
not re-optimized at each scale; Section~\ref{sec:concl} draws the implications for how scaling studies should report
it.

\subsection{A Controlled Experiment}\label{sec:tech:experiment}

The public sweeps confound technology with design, units, and learning-rate schedules. Our controlled experiment
(Section~\ref{sec:data}) holds everything fixed except data quality. Two ``labs'' train the same eight model sizes
(0.4M to 49M non-embedding parameters) on FineWeb-Edu and on FineWeb \citep{penedo2024fineweb}, with a shared
8,192-token tokenizer, token budgets from 25M to 800M, warmup--stable--decay schedules whose cooldown branches give an
annealed endpoint at every budget \citep{hagele2024scaling}, a learning rate calibrated to each width, and common
random numbers. Every model is evaluated on both validation sets, which provides the common output that the
neutrality tests require. Figure~\ref{fig:experiment} shows the fitted technologies of the two labs. Four results
follow.
\textcolor{red}{[TBD-m9: $\sigma^*$ and $\sigma^*_\kappa$ by lab, with $\kappa$ free, on a common output; compare with the 0.7--0.8 range above.]}
\textcolor{red}{[TBD-m9: neutrality of data quality: which of Hicks-neutral, data-augmenting, parameter-augmenting, or asymptote-only is rejected; implied $\chi$ and its effect on $M^*$ and $\hat w$; sign relative to Lemma~\ref{lem:alloc} (better data $\Rightarrow$ fewer tokens per parameter).]}
\textcolor{red}{[TBD-m9: extrapolation check: technology fitted on $M\le100$ (or 341) and evaluated at $M\approx1{,}000$--$2{,}000$; ratio $w_{\rm extrap}/w_{\rm full}$; whether $\varepsilon_D$ is overstated and the direction of the implied bias in Section~\ref{sec:wedge}.]}
\textcolor{red}{[TBD-m9: seed noise: s.d.\ of $\ln L$ across seeds against the residual s.d.\ of the fitted technology.]}

\begin{figure}[tp]
\centering
\fbox{\parbox{0.9\textwidth}{\centering\vspace{2em}\textcolor{red}{[TBD-m9: two panels. (a) Loss against $D$
by width for both labs on a common validation set, with fitted curves. (b) Isoquants and expansion paths of both labs.
File: output/figures/m9\_*.]}\vspace{2em}}}
\caption{A Controlled Experiment: Two Labs That Differ in Data Quality}\label{fig:experiment}
\begin{figurenotes}
\textcolor{red}{[TBD-m9: design, output, estimator, and bootstrap description.]}
\end{figurenotes}
\end{figure}
